% Part: normal-modal-logic
% Chapter: tableaux
% Section: countermodels

\documentclass[../../../include/open-logic-section]{subfiles}

\begin{document}

\olfileid{nml}{tab}{cou}
      
\olsection{النماذج المضادة من \usetoken{P}{tableau}}

\begin{explain}
  لا تقتصر فائدة برهان التمام على أن~$\Entails !A$ يستلزم
  $\Proves !A$؛ بل منه نستخرج نموذجًا مضادًا لـ~$!A$
  متى كان~$\Entails/ !A$. وهذه الطريقة في~$\Log{K}$ \emph{إجراء قرار}.
  فلنفرض، مثلًا،~$\Entails/ !A$. يعطينا برهان
  \olref[cpl]{prop:complete-tableau} طريقًا لبناء !!{tableau} تام،
  وهذا البناء ينتهي. فإن القواعد القضية في~$\Log{K}$ لا تزيد إلا
  !!{formula}s ذات بادئات أقل تعقيدًا؛ فلا حاجة إلى إعادة قاعدة
  قضية في الفرع على الصيغة الموقعة نفسها~$\sFmla{S}{!A}[\sigma]$.
  وأما إنشاء البادئات الجديدة فمختص بالقاعدة
  \iftag{prvBox}{$\TRule{\False}{\Box}$}{}\iftag{notprvBox,notprvDiamond}{}{
    و}\iftag{prvDiamond}{$\TRule{\True}{\Diamond}$}{}
  \iftag{notprvBox,notprvDiamond}{، ولا يلزم تطبيقها}{، ولا يلزم تطبيقهما}
  على الصيغة الموقعة الواحدة إلا مرةً واحدةً أيضًا (وتنتج من كل تطبيق
  بادئةٌ جديدةٌ واحدةٌ).
  وقد يلزم تطبيق
  \iftag{prvBox}{$\TRule{\True}{\Box}$}{}\iftag{notprvBox,notprvDiamond}{}{
    و}\iftag{prvDiamond}{$\TRule{\False}{\Diamond}$}{}
  \iftag{notprvBox,notprvDiamond}{عدة مرات}{عدة مرات}، ولكن مرة واحدة
  فقط لكل بادئة. وبيان تناهي البادئات أن كل بادئة جديدة تحمل صيغةً دون
  الصيغة المولّدة لها عند بادئتها الأم في العمق الجهوي؛ فلا تطول سلسلة
  البادئات على أكبر عمق جهوي في الصيغ الابتدائية. ولأن الصيغ الجزئية
  الموقعة الممكنة عند كل بادئة منتهية، تتناهى البادئات الجديدة. ولذلك
  ينتهي البناء بعد عدد منتهٍ من المراحل إما إلى فرع مغلق وإما إلى فرع تام.
  % تصحيح برهاني (0469-I1): قُيّد تناهي البادئات صراحةً بالعمق الجهوي المتناقص.

  فإذا حصل جدول دلالي فيه فرع مفتوح تام، أمكن ببرهان
  \olref[cpl]{thm:tableau-completeness} تعيين نموذج يرضي المجموعة
  الأصلية من ال!!{formula}s ذات البادئات. فليست النتيجة مجرد أن
  $\Gamma \Entails !A$ يقتضي وجود !!{tableau} مغلق، ومن ثم
  $\Gamma \Proves !A$. بل إذا طلبنا ال!!{tableau} المغلق بهذه
  الطريقة، ثم انتهينا إلى !!{tableau} «تام» لا يزال فيه فرع مفتوح،
  عرفنا~$\Gamma \Entails/ !A$ واستخرجنا نموذجًا مضادًا بالفعل.
\end{explain}

\iftag{prvBox}{
\begin{ex}
  نعلم أن $\Proves/ \Box(p \lor q) \lif (\Box p \lor \Box q)$.
  ولنشرع في بناء جدولها الدلالي:
  \begin{oltableau}
    [\pFmla{\False}{\Box(p \lor q) \lif (\Box p \lor \Box q)}{1},
      just = \TAss, checked
      [\pFmla{\True}{\Box(p \lor q)}{1},
        just = {\TRule{\False}{\lif}[1]}, 
        [\pFmla{\False}{\Box p \lor \Box q}{1},
          just = {\TRule{\False}{\lif}[1]}, checked
          [\pFmla{\False}{\Box p}{1},
            just = {\TRule{\False}{\lor}[3]}, checked
            [\pFmla{\False}{\Box q}{1},
              just = {\TRule{\False}{\lor}[3]}, checked
              [\pFmla{\False}{p}{1.1}, 
                just = {\TRule{\False}{\Box}[4]}, checked
                [\pFmla{\False}{q}{1.2}, 
                  just = {\TRule{\False}{\Box}[5]}, checked
                ]
              ]
            ]
          ]
        ]
      ]
    ]
  \end{oltableau}
  ولم يتم ال!!{tableau} بعد. فلننظر في السطر الذي بقي بلا علامة
  فحص، وهو ال!!{formula} ذات البادئة
  $\sFmla{\True}{\Box(p \lor q)}[1]$ في السطر~$2$.
  يستلزم تمام الجدول إجراء~$\TRule{\True}{\Box}$ عند كل بادئة
  صالحة مستعملة في الفرع، وهما هنا~$1.1$ و$1.2$:
  \begin{oltableau}
    [\pFmla{\False}{\Box(p \lor q) \lif (\Box p \lor \Box q)}{1},
      just = \TAss, checked
      [\pFmla{\True}{\Box(p \lor q)}{1},
        just = {\TRule{\False}{\lif}[1]}, 
        [\pFmla{\False}{\Box p \lor \Box q}{1},
          just = {\TRule{\False}{\lif}[1]}, checked
          [\pFmla{\False}{\Box p}{1},
            just = {\TRule{\False}{\lor}[3]}, checked
            [\pFmla{\False}{\Box q}{1},
              just = {\TRule{\False}{\lor}[3]}, checked
              [\pFmla{\False}{p}{1.1}, 
                just = {\TRule{\False}{\Box}[4]}, checked
                [\pFmla{\False}{q}{1.2}, 
                  just = {\TRule{\False}{\Box}[5]}, checked
                  [\pFmla{\True}{p \lor q}{1.1}, 
                    just = {\TRule{\True}{\Box}[2]}
                    [\pFmla{\True}{p \lor q}{1.2}, 
                      just = {\TRule{\True}{\Box}[2]}
                    ]
                  ]
                ]
              ]
            ]
          ]
        ]
      ]
    ]
  \end{oltableau}
  فبقيت الأسطر 2 و8 و9 بلا علامات فحص. ولم تحدث بادئة جديدة،
  فنُجري~$\TRule{\True}{\lor}$ على السطرين~8 و~9 في كل فرع
  ينتج ما لم ينغلق:
  \begin{oltableau}
    [\pFmla{\False}{\Box(p \lor q) \lif (\Box p \lor \Box q)}{1},
      just = \TAss, checked
      [\pFmla{\True}{\Box(p \lor q)}{1},
        just = {\TRule{\False}{\lif}[1]}, checked
        [\pFmla{\False}{\Box p \lor \Box q}{1},
          just = {\TRule{\False}{\lif}[1]}, checked
          [\pFmla{\False}{\Box p}{1},
            just = {\TRule{\False}{\lor}[3]}, checked
            [\pFmla{\False}{\Box q}{1},
              just = {\TRule{\False}{\lor}[3]}, checked
              [\pFmla{\False}{p}{1.1}, 
                just = {\TRule{\False}{\Box}[4]}, checked
                [\pFmla{\False}{q}{1.2}, 
                  just = {\TRule{\False}{\Box}[5]}, checked
                  [\pFmla{\True}{p \lor q}{1.1}, 
                    just = {\TRule{\True}{\Box}[2]}, checked
                    [\pFmla{\True}{p \lor q}{1.2}, 
                      just = {\TRule{\True}{\Box}[2]}, checked
                      [\pFmla{\True}{p}{1.1},
                        just = {\TRule{\True}{\lor}[8]}, checked, close
                      ]
                      [\pFmla{\True}{q}{1.1},
                        just = {\TRule{\True}{\lor}[8]}, checked
                        [\pFmla{\True}{p}{1.2},
                          just = {\TRule{\True}{\lor}[9]}, checked]
                        [\pFmla{\True}{q}{1.2},
                          just = {\TRule{\True}{\lor}[9]}, checked, close]
                      ]
                    ]
                  ]
                ]
              ]
            ]
          ]
        ]
      ]
    ]
  \end{oltableau}
  لم يبق إلا فرع واحد مفتوح، وهو تام. فنأخذ منه عوالم النموذج
  $W = \{1, 1.1, 1.2\}$، وهي جميع البادئات في ذلك الفرع،
  ونجعل علاقة الإتاحة
  $R = \{\tuple{1, 1.1}, \tuple{1, 1.2}\}$.
  وأما الإسناد فـ~$V(p) =
  \{1.2\}$، لورود~$\sFmla{\True}{p}[1.2]$ في السطر~11؛
  و$V(q) = \{1.1\}$، لورود
  $\sFmla{\True}{q}[1.1]$ في السطر~10. وهذا هو النموذج المرسوم في
  \olref{fig:counter-Box}. فتحقق من أنه مضاد للصيغة
  $\Box(p \lor q) \lif (\Box p \lor \Box q)$.
  \begin{figure}
  \begin{center}
    \begin{tikzpicture}[modal]
      \node[world] (w1) [label={[align=right]right:\mFalse{p}\\ \mFalse{q}}]{$1$}; 
      \node[world] (w2) [label={[align=right]right:\mFalse{p}\\ \mTrue{q}},
        above left=of w1]{$1.1$}; 
      \node[world] (w3) [label={[align=right]right:\mTrue{p}\\ \mFalse{q}},
        above right=of w1] {$1.2$};
      \draw[->] (w1) to (w2);
      \draw[->] (w1) to (w3);
    \end{tikzpicture}
  \end{center}
  \caption{نموذج مضاد لـ$\Box(p \lor q) \lif (\Box p \lor \Box q)$.}
  \ollabel{fig:counter-Box}
\end{figure}
\end{ex}          
}
{
\begin{ex}
  نعلم أن $\Proves/ (\Diamond p \land \Diamond q) \lif \Diamond(p
  \land q)$. ولنبدأ جدولها الدلالي بما يأتي:
  \begin{oltableau}
    [\pFmla{\False}{(\Diamond p \land \Diamond q) \lif \Diamond(p \land q)}{1},
      just = \TAss, checked
      [\pFmla{\True}{\Diamond p \land \Diamond q}{1},
        just = {\TRule{\False}{\lif}[1]}, checked
        [\pFmla{\False}{\Diamond(p \land q)}{1},
          just = {\TRule{\False}{\lif}[1]}
          [\pFmla{\True}{\Diamond p}{1},
            just = {\TRule{\True}{\land}[2]}, checked
            [\pFmla{\True}{\Diamond q}{1},
              just = {\TRule{\True}{\land}[2]}, checked
              [\pFmla{\True}{p}{1.1}, 
                just = {\TRule{\True}{\Diamond}[4]}, checked
                [\pFmla{\True}{q}{1.2}, 
                  just = {\TRule{\True}{\Diamond}[5]}, checked
                ]
              ]
            ]
          ]
        ]
      ]
    ]
  \end{oltableau}
  ولم يتم ال!!{tableau} بعد؛ إذ بقي سطر واحد بلا علامة فحص، فيه
  ال!!{formula} ذات البادئة
  $\sFmla{\False}{\Diamond(p \land q)}[1]$، وهو السطر~$3$.
  ولإتمام الجدول نجري~$\TRule{\False}{\Diamond}$ عند كل بادئة
  صالحة مستعملة في الفرع، أي~$1.1$ و$1.2$:
  \begin{oltableau}
    [\pFmla{\False}{(\Diamond p \land \Diamond q) \lif \Diamond(p \land q)}{1},
      just = \TAss, checked
      [\pFmla{\True}{\Diamond p \land \Diamond q}{1},
        just = {\TRule{\False}{\lif}[1]}, checked
        [\pFmla{\False}{\Diamond (p \land q)}{1},
          just = {\TRule{\False}{\lif}[1]}
          [\pFmla{\True}{\Diamond p}{1},
            just = {\TRule{\True}{\land}[2]}, checked
            [\pFmla{\True}{\Diamond q}{1},
              just = {\TRule{\True}{\land}[2]}, checked
              [\pFmla{\True}{p}{1.1}, 
                just = {\TRule{\True}{\Diamond}[4]}, checked
                [\pFmla{\True}{q}{1.2}, 
                  just = {\TRule{\True}{\Diamond}[5]}, checked
                  [\pFmla{\False}{p \land q}{1.1}, 
                    just = {\TRule{\False}{\Diamond}[3]}
                    [\pFmla{\False}{p \land q}{1.2}, 
                      just = {\TRule{\False}{\Diamond}[3]}
                    ]
                  ]
                ]
              ]
            ]
          ]
        ]
      ]
    ]
  \end{oltableau}
  وتبقى الآن الأسطر 3 و8 و9 بلا علامات فحص. ولما لم تُستحدث بادئة،
  نُجري~$\TRule{\False}{\land}$ على السطرين~8 و~9 في كل فرع
  ناتج لا يزال مفتوحًا:
  \begin{oltableau}
    [\pFmla{\False}{(\Diamond p \land \Diamond q) \lif \Diamond(p \land q)}{1},
      just = \TAss, checked
      [\pFmla{\True}{\Diamond p \land \Diamond q}{1},
        just = {\TRule{\False}{\lif}[1]}, checked
        [\pFmla{\False}{\Diamond(p \land q)}{1},
          just = {\TRule{\False}{\lif}[1]}
          [\pFmla{\True}{\Diamond p}{1},
            just = {\TRule{\True}{\land}[2]}, checked
            [\pFmla{\True}{\Diamond q}{1},
              just = {\TRule{\True}{\land}[2]}, checked
              [\pFmla{\True}{p}{1.1}, 
                just = {\TRule{\True}{\Diamond}[4]}, checked
                [\pFmla{\True}{q}{1.2}, 
                  just = {\TRule{\True}{\Diamond}[5]}, checked
                  [\pFmla{\False}{p \land q}{1.1}, 
                    just = {\TRule{\False}{\Diamond}[3]}, checked
                    [\pFmla{\False}{p \land q}{1.2}, 
                      just = {\TRule{\False}{\Diamond}[3]}, checked
                      [\pFmla{\False}{p}{1.1},
                        just = {\TRule{\False}{\land}[8]}, checked, close
                      ]
                      [\pFmla{\False}{q}{1.1},
                        just = {\TRule{\False}{\land}[8]}, checked
                        [\pFmla{\False}{p}{1.2},
                          just = {\TRule{\False}{\land}[9]}, checked]
                        [\pFmla{\False}{q}{1.2},
                          just = {\TRule{\False}{\land}[9]}, checked, close]
                      ]
                    ]
                  ]
                ]
              ]
            ]
          ]
        ]
      ]
    ]
  \end{oltableau}
  فالفرع المفتوح الباقي واحد، وقد تم. نجعل عوالم النموذج
  $W = \{1, 1.1, 1.2\}$، إذ لا بادئات سواها في هذا الفرع،
  وعلاقة الإتاحة
  $R = \{\tuple{1, 1.1}, \tuple{1, 1.2}\}$.
  ونضع~$V(p) =
  \{1.1\}$ لورود~$\sFmla{\True}{p}[1.1]$ في السطر~6،
  و$V(q) = \{1.2\}$ لورود
  $\sFmla{\True}{q}[1.2]$ في السطر~7. والنموذج في
  \olref{fig:counter-Diamond}؛ فتحقق من أنه مضاد للصيغة
  $(\Diamond p \land \Diamond q) \lif \Diamond (p \land q)$.
  \begin{figure}
  \begin{center}
    \begin{tikzpicture}[modal]
      \node[world] (w1) [label={[align=right]right:\mFalse{p}\\ \mFalse{q}}]{$1$}; 
      \node[world] (w2) [label={[align=right]right:\mTrue{p}\\ \mFalse{q}},
        above left=of w1]{$1.1$}; 
      \node[world] (w3) [label={[align=right]right:\mFalse{p}\\ \mTrue{q}},
        above right=of w1] {$1.2$};
      \draw[->] (w1) to (w2);
      \draw[->] (w1) to (w3);
    \end{tikzpicture}
  \end{center}
  \caption{نموذج مضاد لـ$(\Diamond p \land \Diamond q) \lif
    \Diamond (p \land q)$.}
  \ollabel{fig:counter-Diamond}
\end{figure}
\end{ex}          
}

\end{document}
